% GENERATED FILE. Edit canonical lessons, then run npm run generate:books.
% canonical-source-hash: fnv1a64:60b72534cacaf057
% canonical-lessons: TE-C332-appatike, TE-C332-antalo, TE-C332-hathattuga, TE-C332-ksanam, TE-C332-arucu

\chapter{Already, Meanwhile, Suddenly}
\label{ch:te-already-meanwhile-suddenly}
\hlchaptermodality{332}

\begin{chapteropening}
\textbf{By the end of this chapter:} \emph{I can say already by then, in the meantime, all of a sudden, a moment and to yell.}

Complete the last lesson of chapter 332: I can say already by then, in the meantime, all of a sudden, a moment and to yell.
\end{chapteropening}

\section[\texorpdfstring{\te{అప్పటికే} (appaṭikē)}{అప్పటికే (appaṭikē)}]{\te{అప్పటికే} (appaṭikē) — already by then}

\label{lesson:TE-C332-appatike}



\begin{quote}
Before the new one: say the Telugu for to rescue, then the Telugu for to run for it.
\end{quote}

\subsection*{You'll want to know: \te{అప్పటికే}}
\textbf{\te{అప్పటికే}} — \emph{appaṭikē} — \textquotedblleft{}already by then\textquotedblright{}.

\begin{quote}
\te{పోలీసులు} \te{వచ్చారు}, \te{కానీ} \te{దొంగ} \te{అప్పటికే} \te{పారిపోయాడు}.
\end{quote}

\begin{quote}
\emph{pōlīsulu vaccāru, kānī doṅga appaṭikē pāripōyāḍu.}
\end{quote}

\begin{quote}
The police came, but the thief had already run away.
\end{quote}

\te{అప్పటికి}, by then, with -\te{ఏ}, the ending that stresses. Compare \te{ఇప్పటికే}, already by now.

\begin{cousinweb}
\textbf{\te{అప్పటికే}} (\emph{appaṭikē}) is native Telugu: \emph{appaṭiki}, \textquotedblleft{}by then\textquotedblright{}, from \emph{appuḍu}, \textquotedblleft{}at that time\textquotedblright{} (DEDR 4559), with the stressing ending \emph{-ē}.
\end{cousinweb}

\begin{grammarlens}[title={What you've built}]
One more word for telling what happened.
\end{grammarlens}

\subsection*{Your turn}
\begin{itemize}
\raggedright
  \item \emph{Say it:} \emph{appaṭikē}
  \item \emph{Say it:} \emph{appaṭikē}, once more
  \item \emph{Say it:} \emph{pōlīsulu vaccāru, kānī doṅga appaṭikē pāripōyāḍu.} — The police came, but the thief had already run away.
  \item \emph{Recall:} say the Telugu for to rescue, then the Telugu for to run for it, then say \emph{appaṭikē} again
\end{itemize}

\begin{tcolorbox}[breakable,colback=teal!4,colframe=teal!35!black,title={Before you move on}]
What does \te{అప్పటికే} mean? (\textquotedblleft{}Already by then\textquotedblright{}.) Say it once more.
\end{tcolorbox}
\section[\texorpdfstring{\te{అంతలో} (antalō)}{అంతలో (antalō)}]{\te{అంతలో} (antalō) — in the meantime}

\label{lesson:TE-C332-antalo}



\begin{quote}
Before the new one: say the Telugu for to run for it, then the Telugu for already by then.
\end{quote}

\subsection*{You'll want to know: \te{అంతలో}}
\textbf{\te{అంతలో}} — \emph{antalō} — \textquotedblleft{}in the meantime\textquotedblright{}.

\begin{quote}
\te{నేను} \te{ఫోను} \te{చేశాను}, \te{అంతలో} \te{అతనే} \te{వచ్చాడు}.
\end{quote}

\begin{quote}
\emph{nēnu phōnu cēśānu, antalō atanē vaccāḍu.}
\end{quote}

\begin{quote}
I phoned, and in the meantime he came himself.
\end{quote}

\te{అంత}, that much, with -\te{లో}, in: within that span. It joins two happenings that overlap.

\begin{cousinweb}
\textbf{\te{అంతలో}} (\emph{antalō}) is native Telugu: \emph{anta}, \textquotedblleft{}that much\textquotedblright{}, from the far demonstrative \emph{a-} (DEDR 1), with the locative ending \emph{-lō}, \textquotedblleft{}in\textquotedblright{}.
\end{cousinweb}

\begin{grammarlens}[title={What you've built}]
One more word for telling what happened.
\end{grammarlens}

\subsection*{Your turn}
\begin{itemize}
\raggedright
  \item \emph{Say it:} \emph{antalō}
  \item \emph{Say it:} \emph{antalō}, once more
  \item \emph{Say it:} \emph{nēnu phōnu cēśānu, antalō atanē vaccāḍu.} — I phoned, and in the meantime he came himself.
  \item \emph{Recall:} say the Telugu for to run for it, then the Telugu for already by then, then say \emph{antalō} again
\end{itemize}

\begin{tcolorbox}[breakable,colback=teal!4,colframe=teal!35!black,title={Before you move on}]
What does \te{అంతలో} mean? (\textquotedblleft{}In the meantime\textquotedblright{}.) Say it once more.
\end{tcolorbox}
\section[\texorpdfstring{\te{హఠాత్తుగా} (haṭhāttugā)}{హఠాత్తుగా (haṭhāttugā)}]{\te{హఠాత్తుగా} (haṭhāttugā) — all of a sudden}

\label{lesson:TE-C332-hathattuga}



\begin{quote}
Before the new one: say the Telugu for already by then, then the Telugu for in the meantime.
\end{quote}

\subsection*{You'll want to know: \te{హఠాత్తుగా}}
\textbf{\te{హఠాత్తుగా}} — \emph{haṭhāttugā} — \textquotedblleft{}all of a sudden\textquotedblright{}.

\begin{quote}
\te{హఠాత్తుగా} \te{కరెంటు} \te{పోయింది}.
\end{quote}

\begin{quote}
\emph{haṭhāttugā kareṇṭu pōyindi.}
\end{quote}

\begin{quote}
All of a sudden the power went.
\end{quote}

Like \te{అకస్మాత్తుగా}, it ends in -\te{గా}, the ending that makes an adverb.

\begin{cousinweb}
\textbf{\te{హఠాత్తుగా}} (\emph{haṭhāttugā}) is built on Sanskrit \emph{haṭhāt}, \textquotedblleft{}by force, forcibly\textquotedblright{}, a form of \emph{haṭha}, \textquotedblleft{}force\textquotedblright{} (Monier-Williams), with the Telugu adverb ending \emph{-gā}.
\end{cousinweb}

\begin{grammarlens}[title={What you've built}]
One more word for telling what happened.
\end{grammarlens}

\subsection*{Your turn}
\begin{itemize}
\raggedright
  \item \emph{Say it:} \emph{haṭhāttugā}
  \item \emph{Say it:} \emph{haṭhāttugā}, once more
  \item \emph{Say it:} \emph{haṭhāttugā kareṇṭu pōyindi.} — All of a sudden the power went.
  \item \emph{Recall:} say the Telugu for already by then, then the Telugu for in the meantime, then say \emph{haṭhāttugā} again
\end{itemize}

\begin{tcolorbox}[breakable,colback=teal!4,colframe=teal!35!black,title={Before you move on}]
What does \te{హఠాత్తుగా} mean? (\textquotedblleft{}All of a sudden\textquotedblright{}.) Say it once more.
\end{tcolorbox}
\section[\texorpdfstring{\te{క్షణం} (kṣaṇaṁ)}{క్షణం (kṣaṇaṁ)}]{\te{క్షణం} (kṣaṇaṁ) — a moment; an instant}

\label{lesson:TE-C332-ksanam}



\begin{quote}
Before the new one: say the Telugu for in the meantime, then the Telugu for all of a sudden.
\end{quote}

\subsection*{You'll want to know: \te{క్షణం}}
\textbf{\te{క్షణం}} — \emph{kṣaṇaṁ} — \textquotedblleft{}a moment; an instant\textquotedblright{}.

\begin{quote}
\te{ఒక్క} \te{క్షణంలో} \te{అంతా} \te{మారిపోయింది}.
\end{quote}

\begin{quote}
\emph{okka kṣaṇaṁlō antā māripōyindi.}
\end{quote}

\begin{quote}
In a single moment everything changed.
\end{quote}

Plural: \te{క్షణాలు}, \emph{kṣaṇālu}.

\te{ఒక్క} \te{క్షణం}, one moment, also asks someone to wait.

\begin{cousinweb}
\textbf{\te{క్షణం}} (\emph{kṣaṇaṁ}) is borrowed from Sanskrit \emph{kṣaṇa}, \textquotedblleft{}the twinkling of an eye, a moment\textquotedblright{} (CDIAL 3642).
\end{cousinweb}

\begin{grammarlens}[title={What you've built}]
One more word for telling what happened.
\end{grammarlens}

\subsection*{Your turn}
\begin{itemize}
\raggedright
  \item \emph{Say it:} \emph{kṣaṇaṁ}
  \item \emph{Say it:} \emph{kṣaṇaṁ}, once more
  \item \emph{Say it:} \emph{okka kṣaṇaṁlō antā māripōyindi.} — In a single moment everything changed.
  \item \emph{Recall:} say the Telugu for in the meantime, then the Telugu for all of a sudden, then say \emph{kṣaṇaṁ} again
\end{itemize}

\begin{tcolorbox}[breakable,colback=teal!4,colframe=teal!35!black,title={Before you move on}]
What does \te{క్షణం} mean? (\textquotedblleft{}A moment; an instant\textquotedblright{}.) Say it once more.
\end{tcolorbox}
\section[\texorpdfstring{\te{అరుచు} (arucu)}{అరుచు (arucu)}]{\te{అరుచు} (arucu) — to yell; to shout}

\label{lesson:TE-C332-arucu}



\begin{quote}
Before the new one: say the Telugu for all of a sudden, then the Telugu for a moment.
\end{quote}

\subsection*{You'll want to know: \te{అరుచు}}
\textbf{\te{అరుచు}} — \emph{arucu} — \textquotedblleft{}to yell; to shout\textquotedblright{}.

\begin{quote}
\te{పామును} \te{చూసి} \te{పిల్లవాడు} \te{గట్టిగా} \te{అరిచాడు}.
\end{quote}

\begin{quote}
\emph{pāmunu cūsi pillavāḍu gaṭṭigā aricāḍu.}
\end{quote}

\begin{quote}
Seeing a snake, the boy yelled loudly.
\end{quote}

Past (I …): \te{అరిచాను}, \emph{aricānu}.

You know its noun, \te{అరుపు}, a shout.

\begin{cousinweb}
\textbf{\te{అరుచు}} (\emph{arucu}) is native Dravidian: DEDR 367 lists Telugu \emph{ārcu}, \emph{ārucu}, \textquotedblleft{}to cry aloud, shout\textquotedblright{}, with Tamil \emph{ār}, \textquotedblleft{}to shout, roar\textquotedblright{}.
\end{cousinweb}

\begin{grammarlens}[title={What you've built}]
One more word for telling what happened.
\end{grammarlens}

\subsection*{Your turn}
\begin{itemize}
\raggedright
  \item \emph{Say it:} \emph{arucu}
  \item \emph{Say it:} \emph{arucu}, once more
  \item \emph{Say it:} \emph{pāmunu cūsi pillavāḍu gaṭṭigā aricāḍu.} — Seeing a snake, the boy yelled loudly.
  \item \emph{Recall:} say the Telugu for all of a sudden, then the Telugu for a moment, then say \emph{arucu} again
\end{itemize}

\begin{tcolorbox}[breakable,colback=teal!4,colframe=teal!35!black,title={Before you move on}]
What does \te{అరుచు} mean? (\textquotedblleft{}To yell; to shout\textquotedblright{}.) Say it once more.
\end{tcolorbox}
